\documentclass{article}

\usepackage[final]{corl_2026}
\usepackage{booktabs}
\usepackage{amsmath}

% Course footer instead of the CoRL conference notice.
\makeatletter
\renewcommand{\@noticestring}{Course Project, 16-831 Introduction to Robot Learning, Fall 2026, CMU}
\makeatother

\title{How Much Should the Robot Help? Adaptive Assistance Levels for Diffusion-Based Shared Autonomy}

\author{
  Jarrod Yang \\
  \texttt{jiayuy2} \\
  \texttt{jiayuy2@andrew.cmu.edu} \\
  \And
  Ziyu Xu \\
  \texttt{rickx} \\
  \texttt{rickx@andrew.cmu.edu} \\
  \And
  Runhan Wang \\
  \texttt{runhanw} \\
  \texttt{runhanw@andrew.cmu.edu} \\
  \And
  Runhui Long \\
  \texttt{runhuil} \\
  \texttt{runhuil@andrew.cmu.edu} \\
}

\begin{document}
\maketitle

%===============================================================================

\begin{abstract}
Wearable interfaces such as surface-EMG wristbands promise intuitive robot-arm control, but their commands are noisy, delayed, and vary across users and over time. Diffusion-based shared autonomy~\citep{yoneda2023noise} corrects user actions with a copilot whose forward-diffusion ratio $\gamma$ trades off correcting noise against preserving the user's intent, but $\gamma$ is fixed for all users and over time. We reproduce this method on Block Pushing and study whether, and how, the assistance level should adapt. We will build simulated users calibrated to real EMG decoding errors and compare a uniform $\gamma$, bandit calibration, and a reinforcement-learning policy that sets $\gamma$ at every step, measuring task success, how often the copilot overrides the user's goal, and the feedback needed to adapt.
\end{abstract}

\keywords{Shared Autonomy, Diffusion Models, Personalization}

%===============================================================================

\section{Introduction}

Imagine controlling a robot arm with a muscle-sensing wristband. On a good day the decoded commands are close to what you intend; after an hour of use, fatigue and a slightly shifted band make them jittery and delayed. A helpful robot should step in more when your commands get worse and step back when you are precise, since too much help can pursue the wrong goal.

Shared autonomy combines user input with robot autonomy~\citep{dragan2013policy, javdani2015shared} and can be learned with deep RL~\citep{reddy2018shared, schaff2020residual}. \citet{yoneda2023noise} showed that a diffusion model~\citep{ho2020ddpm, chi2023diffusion} trained only on expert demonstrations can act as a copilot: the user's action is partially forward-diffused and then denoised conditioned on the observation. The ratio $\gamma\in[0,1]$ is an explicit assistance knob ($\gamma=0$ executes the user's action; $\gamma\to1$ ignores it), but it is fixed per task and shared by all users. Adaptive assistance itself is not new: arbitration can be customized per user~\citep{gopinath2017customized} or co-adapted with the human~\citep{nikolaidis2017mutual}, and IDA~\citep{mcmahan2024ida} lets a diffusion copilot intervene only when its action is better than the human's for every possible goal. IDA, however, is a binary, value-based switch and does not model users whose input quality differs across people and drifts over time, which is the regime of EMG and other wearable interfaces~\citep{losey2020latent}.

Our preliminary results make this concrete. Using the official checkpoints, we reproduced the paper's Block Pushing results (Table~\ref{tab:repro}) and swept 11 simulated users $\times$ 11 values of $\gamma$ (24{,}200 paired episodes). Consistent with the paper's own sweep, the optimal $\gamma$ grows with user noise (0.15 to 0.4), and a large $\gamma$ sharply increases wrong-goal outcomes. Yet a uniform $\gamma=0.25$ reaches 0.748 mean success, only 1.8 points (cross-fitted) below a per-user oracle, and per-user bandits~\citep{auer2002ucb, russo2018ts} did not beat it within 80 episodes. We hypothesize that the surrogate users understate the diversity of real interfaces and that episode-level choices are too coarse.

\begin{table}[h]
  \centering
  \small
  \caption{Reproduced Block Pushing results (correct-goal rate; noise level 0.6, $\gamma=0.2$; ours: 300 episodes per cell, 95\% CI).}
  \label{tab:repro}
  \begin{tabular}{lcccc}
    \toprule
    & \multicolumn{2}{c}{Noisy user} & \multicolumn{2}{c}{Laggy user} \\
    \cmidrule(lr){2-3}\cmidrule(lr){4-5}
    & No assist & Assisted & No assist & Assisted \\
    \midrule
    Paper & 62.3\% & 75.7\% & 42.0\% & 74.3\% \\
    Ours  & 63.3 $\pm$ 5.5\% & 73.3 $\pm$ 5.0\% & 44.7 $\pm$ 5.6\% & 75.7 $\pm$ 4.9\% \\
    \bottomrule
  \end{tabular}
\end{table}

We therefore study \emph{how much the robot should help} under realistic, user-specific, and drifting input noise, by adapting $\gamma$ at every step:
\begin{enumerate}\setlength{\itemsep}{1pt}
  \item \textbf{EMG-like users.} Noise models with decoding latency, gesture confusions, and slow drift (electrode shift, fatigue), fitted to public sEMG data~\citep{atzori2014ninapro, salter2024emg2pose, ctrllabs2025generic} or, if permitted, our lab's data.
  \item \textbf{Adaptive assistance.} A per-step $\gamma$ policy trained with PPO~\citep{schulman2017ppo} or SAC~\citep{haarnoja2018sac} on task-agnostic features (e.g., the user--copilot action discrepancy), rewarded for the user's goal and penalized for the other one, with the copilot kept fixed. Baselines: uniform $\gamma$, a rule that raises $\gamma$ with the discrepancy, IDA-style switching, and bandits with a population prior.
  \item \textbf{Evaluation.} Correct-goal, wrong-goal, and timeout rates over multiple seeds, the feedback needed to adapt, robustness to within-session drift, and ablations on the policy's inputs.
\end{enumerate}
Core experiments run in simulation. As extensions, we will retrain the copilot for pick-and-place in MuJoCo, a task the paper does not cover that adds grasping to the EMG commands, and, as a stretch goal, demonstrate it on a SO-ARM101 arm with LeRobot, with team members as pilots.

%===============================================================================

\section{Plan}

\begin{itemize}\setlength{\itemsep}{1pt}
  \item 09/27/2026 -- 10/05/2026 (Finished): CPU-only pipeline around the official code and checkpoints; reproduced Block Pushing results; 24{,}200-episode $\gamma$ sweep and offline bandit study.
  \item 10/06/2026 -- 10/13/2026 (Ongoing): Proposal; EMG data source; shared interfaces for tasks, copilots, and users.
  \item 10/14/2026 -- 11/20/2026 (ToDo, Phase 1, core): Block Pushing with EMG-like users built from Ninapro DB2 wrist gestures, with a re-run $\gamma$ sweep by 10/31; train the $\gamma$ policy and compare it with all baselines (multiple seeds, drift tests, ablations).
  \item 10/20/2026 -- 11/24/2026 (ToDo, Phase 2, extension, in parallel): Pick-and-place in MuJoCo with scripted demonstrations to two place zones, a retrained copilot, and grasp/release gestures; zero-shot transfer and fine-tuning of the $\gamma$ policy. Stop if the copilot does not help noisy users by 11/07.
  \item 11/15/2026 -- 11/30/2026 (ToDo, Phase 3, stretch): SO-ARM101 demo with LeRobot, only if Phase 2 succeeds.
  \item 11/25/2026 -- 12/11/2026 (ToDo): Oral presentation (12/01, 12/03) and final report (12/11).
\end{itemize}

%===============================================================================

\bibliography{references}

\end{document}
